\section{Глаза: холодные, светящиеся, золотые, пустые, кошачьи, усталые, "слишком живые"}

Глаза — деталь с наибольшей эмоциональной пропускной способностью, что мы уже разбирали в 1.2 на примере Тэдзуки: минимальное изменение формы даёт максимальный сдвиг смысла. Здесь стоит каталогизировать конкретные коды не абстрактно, а на уже знакомых примерах.

Золотые и кошачьи глаза Holo (6.1) работают вместе, не по отдельности: золото сигнализирует древность и нечеловеческую ценность (тот же принцип, что у фиолетово-золотой пары в 6.4), а вертикальный кошачий зрачок добавляет к этому witch/trickster-дистанцию (1.1, "дистанция через знание") — читатель понимает, что перед ним существо, которое видит иначе, буквально раньше, чем метафорически. Пустые глаза Rei (1.4) работают ровно противоположно: минимум блика, минимум направленности взгляда — визуальный эквивалент нулевой точки, у которой пока нет собственного вектора.

Между этими двумя полюсами — холодные глаза (минимум блика при сохранённой геометрии, kuudere-сигнал, 3.4), усталые глаза (тяжёлые веки, тени — код исторической тяжести) и опасный крайний случай — "слишком живые" глаза, где избыточная детализация радужки создаёт uncanny-эффект (0.1), потому что глаз перестаёт быть символом и начинает конкурировать с реализмом на его условиях.

Для Рин последний пункт — прямое предостережение: её глаза, когда бы они ни были спроектированы, не должны стремиться к максимальной детализации радужки. Смысл в её случае, вероятно, ближе к комбинации Holo, чем Rei, — светящийся или нечеловеческий код плюс редкий тёплый блик, а не пустота, — код, который совпадает с её dere-формулой (3.1): холодно-точная база, тепло как редкое исключение, а не постоянный фон.
